\documentclass{scribe-or}
\definecolor{darkgreen}{rgb}{0.0, 0.2, 0.13}
\definecolor{ao}{rgb}{0.0, 0.5, 0.0}
\usepackage{circuitikz}
\usepackage{graphicx}
\usepackage{tikz}
\usetikzlibrary{circuits.logic.US,positioning,calc}
\usepackage{karnaugh-map}



\begin{document}

\جلسه{مدل تورینگ و کلاس‌های \lr{R} و  \lr{RE}  }{تمرین سری هفتم و هشتم و نهم}{ساعت 59:23 روز 5 تیر}

\textbf{لطفا پیش از پاسخ‌دادن به تمرین‌ها به نکات زیر توجه کنید: \\}	
{\scriptsize \begin{itemize}
	\item
تمرین از دو بخش سوالات تحویلی و سوالات تکمیلی تشکیل شده است. توجه کنید که پاسخ‌دادن به سوالات تکمیلی نمره‌ی اضافه‌ای ندارد.
	\item 
	ارسال سوال‌ها به فروم‌های اینترنتی و جست‌و‌جوی پاسخ آن‌ها در اینترنت مجاز نیست.
	\item 
	می‌توانید با یکدیگر در حل سوالات مشورت کنید؛ اما باید اولاً راه‌حلتان را با بیان خودتان بنویسید و ثانیاً نام کسانی که با آن‌ها در حل سوال مشورت کرده‌اید را پیش از پاسختان به سوال ذکر کنید.
	\item 
	در صورتی که در مورد تمرین‌ها سوالی و ابهامی داشتید پیشنهاد می‌شود از دستیاران بپرسید.
	در صورت تشخیص مشابهت در راه‌حل‌ها، با فرض عدم تخلف تصحیح صورت خواهد گرفت اما مستندات بدون اطلاع دانشجو به مراجع ذی‌صلاح جهت بررسی، تصمیم و اقدام ارسال خواهد شد.
	
	\item
	دقت لازم را در نوشتن اثبات‌ها و بیان ادعاها به‌خرج دهید. علی‌الاصول هر ادعایی که در پاسخ به تمرین‌ها می‌آورید باید با اثبات همراه باشد؛ مگر آن‌که آن گزاره‌ی مزبور در طول درس اثبات شده باشد و یا سوال صراحتاً گفته باشد که نیازی به اثبات نیست.
	\item 
	برای مرتبط کردن بخش‌های مختلف یک اثبات، به‌جای استفاده از پیکان‌، از کلمات استفاده کنید. همچنین برای هر منظور از سورها (
	$\forall,\exists$) 
	استفاده نکنید. 
	 پاسختان به سوالات باید همراه با توضیحات کافی باشد که مصحح بتواند راه‌حل شما را متوجه شود. متن کتاب مرجع را الگو قرار دهید و پاسختان را طوری بنویسید که هر کسی بتواند آن را دنبال کند و متوجه شود.
	
	 \item 
	پاسختان را در فایلی با نام شماره دانشجوییتان در سامانه اپلود کنید. فرمت فایل ارسالی باید حتما به‌صورت 
	\lr{\textbf{.pdf}}
	باشد. اگر از پاسختان عکس می‌گیرید در نور مناسب این‌کار را بکنید و توجه کنید که تصویر واضح باشد. فایل ارسالی شما نباید نیاز به چرخاندن (\lr{rotatation}) داشته باشد. توجه کنید که پاسخ‌هایی که موارد قبل در آن رعایت نشده باشند یا ناخوانا و مخدوش باشند تصحیح نخواهند شد.
\end{itemize}}
\pagebreak
\section*{تمرینات تحویلی}
\subsection*{سوال ۱}
\textcolor{blue}{(۴۰+۴۰ نمره)}\\
برای هر یک از زبان‌های زیر یک ماشین تورینگ طراحی کنید که آن زبان را تشخیص دهد:
\begin{itemize}
\item 
$L_1 = \{a^nb^m | 0 \leq n \leq m\}$
\item 
$L_2 = \{w \in \{0 , 1\}^* | w = w^R\}$
\end{itemize}

\subsection*{سوال ۲}
\textcolor{blue}{(۷۰ نمره)}\\
در تعریف ماشین تورینگ، نشانگر%
\LTRfootnote{Head}
پس از هر انتقال حالت یا باید به راست برود یا به چپ. حال فرض کنید این قابلیت را به ماشین تورینگ اضافه کنیم که مجاز باشد نشانگرش را لزوما تکان ندهد. یعنی تابع انتقال حالت ماشین جدیدی که ارائه میدهیم به شکل زیر باشد
$$\delta: \ Q\times \Gamma \rightarrow Q\times \Gamma \times \{R,L,S\}$$
که در آن $S$ نماینده تکان ندادنِ نشانگر است. نشان دهید برای هر ماشین $M$ که این قابلیت را دارد، ماشین تورینگ عادی $M'$ داریم که $L(M)=L(M')$ باشد. روش ساخت را کامل و فرمال بنویسید.

\subsection*{سوال ۳}
\textcolor{blue}{(۴۰ نمره)}\\
فرض کنید 
$\Sigma=\{0,\;1,\;...,\;9\}$
و $ A$ زبان تمام رشته‌هایی روی الفبای $\Sigma$ باشد که اگر اعضای آن را به عنوان عدد ده‌دهی در نظر بگیریم، یک عدد اول باشند. به عنوان مثال ۲ و ۰۰۱۱ اعضای $A$ هستند ولی ۲۳۶ عضو $A$ نیست. حال تمام زبان‌هایی مثل $B$ را پیدا کنید که روی الفبای $\Sigma$ تعریف شده‌اند و ضمناً داریم :
$$A \leq_m B$$
\subsection*{سوال ۴}
\textcolor{blue}{(30+30+30+30 نمره)}\\
تصمیم‌پذیری زبان‌های زیر را بررسی کنید. برای هر یک، اگر تصمیم‌پذیر هستند یک ماشین تورینگ تصمیم‌گیر ارائه کنید و اگر تصمیم‌پذیر نیستند، برای ادعایتان برهان بیاورید.
\begin{enumerate}
\item 
$$L_1=\{\left\langle M \right\rangle ~:~ \text{
$M$
یک ماشین تورینگ است و 
$|L(M)|<\infty$}\}$$
\item 
$$L_3=\{\left\langle M , M^\prime\right\rangle ~:~\text{
	$M$
	و 
	$M^\prime$
	دو ماشین تورینگ هستند و 
	$L(M)\cap L(M^\prime)$
نامتناهی است}\}$$
\item 
$$L_4=\{\left\langle G , F\right\rangle ~:~\text{
	$G$
	یک گرامر مستقل از متن و $F$ یک اتوماتای متناهی است و  
	$L(G)\subseteq L(F)$}\}$$
\item 
$$L_5=\{\left\langle M \right\rangle ~:~ \text{
	$M$
	یک ماشین تورینگ است و رشته‌ای وجود دارد که 
	$M$
	در حداکثر ۱۰ انتقال آن را می‌پذیرد}
	\}$$ 
\end{enumerate}
\subsection*{سوال ۵}
\textcolor{blue}{(30+30 نمره)}\\
تشخیص‌پذیری زبان‌های زیر را بررسی کنید. برای هر یک، اگر تشخیص‌پذیر هستند یک ماشین تورینگ تشخیص‌‌دهنده ارائه کنید و اگر تشخیص‌پذیر نیستند، برای ادعایتان برهان بیاورید.
\begin{enumerate}
\item 
$$L_1=\{\left\langle M \right\rangle ~:~ \text{
	$M$
	یک ماشین تورینگ است و 
	$|L(M)|=2$}\}$$
\item 
$$L_2=\{\left\langle M , M^\prime\right\rangle ~:~\text{
	$M$
	و 
	$M^\prime$
	دو ماشین تورینگ هستند و 
	$|L(M)\cap L(M^\prime)|>2$}\}$$
\end{enumerate}
\subsection*{سوال ۶}
\textcolor{blue}{(۴۰+۴۰ نمره)}\\
بررسی کنید که آیا کلاس‌های زبان‌های تشخیص‌پذیر و تصمیم‌پذیر (به ترتیب \lr{RE} و \lr{R}) تحت هومومورفیزم بسته هستند یا خیر.
\section*{تمرینات تکمیلی}
\subsection*{سوال 1}
نشان دهید برای هر دو زبان $A$ و $B$ وجود دارد زبان $C$ که $A\leq_mC$ و $B\leq_mC$ باشد.
\subsection*{سوال 2}
ماشین تورینگی حداقل ۳ نواره طراحی کنید که در دو نوار خود دو عدد (در نمایش باینری) را به عنوان ورودی بگیرد و در نوار سوم ضرب این دو عدد را چاپ کند. 
\subsection*{ سوال ۳}
یک ماشین دو-پشته‌ای، همانند ماشین پشته‌ای است‌ (\lr{PDA}) با این تفاوت که دارای دو پشته است (در واقع تابع انتقال حالت آن به صورت 
$\delta : Q\times \Sigma \times \Gamma \times \Gamma \rightarrow 2^{Q\times\Gamma \times \Gamma}$
است که برد این تابع انتقال حالت، زیرمجموعه‌ای از مجموعه‌ی 
$Q\times\Gamma \times \Gamma$
است). نشان دهید به ازای هر ماشین تورینگ تک‌نواره از دو طرف نامحدود مانند $M$، یک ماشین دو-پشته‌ای $P$ وجود دارد که زبانش با زبان $M$ یکسان باشد. ماشین $P$ را به طور دقیق با ارائه‌ی تابع انتقال حالت توصیف کنید.
\subsection*{ سوال ۴}
در تعریف ماشین تورینگ $k$نواره، هر کدام از نشانگرها پس از هر انتقال حالت یا باید به راست برود یا به چپ. یعنی تابع انتقال حالت آن به شکل زیر است
$$\delta: \ Q\times \Gamma^k \rightarrow Q\times \Gamma^k \times \{R,L\}^k$$
حال فرض کنید این قابلیت را به ماشین تورینگ اضافه کنیم که مجاز باشد نشانگرش را لزوما تکان ندهد. یعنی تابع انتقال حالت ماشین جدیدی که ارائه میدهیم به شکل زیر باشد
$$\delta: \ Q\times \Gamma^k \rightarrow Q\times \Gamma^k \times \{R,L,S\}^k$$
که در آن $S$ نماینده تکان ندادنِ نشانگر است. نشان دهید برای هر ماشین $M$ که این قابلیت را دارد، ماشین تورینگ $k$نواره عادی $M'$ داریم که $L(M)=L(M')$ باشد.
\subsection*{ سوال ۵}
تشخیص‌پذیری و تصمیم‌پذیریِ زبان‌های زیر را بررسی کنید.
\begin{enumerate}
\item 
ورودی: $\langle M,w,k\rangle$\\
خروجی: آیا اگر ماشین تورینگ $M$ را روی $w$ اجرا کنیم، حداکثر از $k$ خانه از حافظه‌اش استفاده می‌کند؟
\item 
ورودی: $\langle M,w\rangle$\\
خروجی: آیا حافظه مصرفی $M$ هنگام اجرا روی $w$ متناهی است؟ به عبارتی $k$ وجود دارد که اگر $M$ را روی $w$ اجرا کنیم حداکثر از $k$ خانه از حافظه‌اش استفاده کند؟
\item 
ورودی: $\langle M\rangle$\\
خروجی: آیا ماشین تورینگ $M$ حین اجرا روی $\epsilon$، هیچگاه نشانگرش را ۵ بار متوالی به سمت راست حرکت می‌دهد؟
\item 
ورودی: عدد $k$\\
خروجی: آیا تعداد رشته‌هایی که ماشین تورینگ شماره ١٠٠٠ روی آن ها توقف نمی‌کند کمتر از $k$ است؟\\[1em]
پی‌نوشت: همانطور که می‌دانید تعداد ماشین‌های تورینگ شمارا است و می‌توان آن‌ها را شماره داد تا ماشین ۱۰۰۰ معنا داشته باشد. برای رفع سو تفاهم، فرض می‌کنیم یک شماره دهی از قبل انجام داده‌ایم و الان برای همه، ماشین شماره ۱۰۰۰ یک ماشین مشخص است و همه می‌دانیم چه ماشینی است.
\item 
ورودی: رشته $w$\\
خروجی: آیا اگر ماشین تورینگ شماره ١٠٠٠ را روی کد خودش اجرا کنیم، متوقف می‌شود؟
\item
ورودی: $\langle M\rangle$\\
خروجی: آیا ماشین تورینگ $M$ حین اجرا روی $\epsilon$، هیچگاه نشانگرش را به سمت چپ حرکت می‌دهد؟
\item 

منظور از $RO$، ماشین‌های تورینگِ تکنواره‌ای است که نوارشان فقط از یک طرف نامتناهی است و نیز قادر به بازنویسیِ بخشی از نوار که روی آن ورودی نوشته شده است، نمی‌باشند. به عنوان تمرین خوب است به این مسئله فکر کنید که زبانِ چنین ماشین‌های تورینگی، حتما یک زبانِ منظم است.
\end{enumerate}

\subsection*{ سوال ۶}
آیا الفبای 
$\Sigma$
و زبان‌های 
$\mathcal{L},\mathcal{L}^\prime \subseteq \Sigma^*$
وجود دارند به‌طوری‌که تحویل چند به یکی\LTRfootnote{many-to-one reduction} بین این دو زبان وجود نداشته باشد؟

\subsection*{ سوال ۷}
ماشین تورینگی طراحی کنید به طوری که هرگاه در ورودی آن رشته‌ی 
$\#1^n\#$
بنویسیم پس از پایان اجرا روی نوار خود رشته‌ی 
$\#1^{\lceil\log_2(n)\rceil}\#$
را چاپ کند.
\end{document}